\begin{terminologyList}
    \term[AllGather]{коллективная операция, при которой каждый участник передает свой фрагмент данных всем остальным участникам, после чего у каждого участника оказывается полный набор фрагментов}
    \term[AllReduce]{коллективная операция, при которой значения от всех участников объединяются с помощью заданной операции, например суммирования, и результат становится доступен всем участникам}
    \term[All-to-All]{коллективная операция обмена данными, при которой каждый участник вычисления передает данные всем остальным участникам и получает данные от них}
    \term[Combine]{этап MoE слоя, на котором результаты вычислений выбранных экспертов собираются обратно и приводятся к порядку исходных входных активаций}
    \term[DeepEP]{специализированная библиотека для высокопроизводительных коммуникаций в MoE моделях, оптимизирующая операции dispatch и combine при экспертном параллелизме}
    \term[Dispatch]{этап MoE слоя, на котором входные активации или токены отправляются к экспертам, выбранным роутером}
    \term[Expert parallelism (EP)]{способ распределенного выполнения MoE моделей, при котором разные эксперты размещаются на разных вычислительных устройствах или узлах}
    \term[Mixture-of-Experts (MoE)]{архитектура нейронной сети, в которой часть полносвязных слоев заменяется набором специализированных экспертных блоков, а роутер выбирает подмножество экспертов для обработки каждого входа}
    \term[NCCL (NVIDIA Collective Communication Library)]{библиотека NVIDIA для коллективных коммуникаций и адресных передач между GPU, широко используемая в распределенном обучении и инференсе нейронных сетей}
    \term[NIXL (NVIDIA Inference Xfer Library)]{библиотека для низкоуровневой высокопроизводительной передачи данных, используемая в работе как одна из реализаций интерфейса передачи тензоров}
    \term[NVSHMEM (NVIDIA OpenSHMEM)]{библиотека, реализующая модель OpenSHMEM для кластеров GPU и предоставляющая односторонний доступ к памяти, синхронизацию и коллективные операции}
    \term[MoonCake Transfer Engine]{библиотека для высокопроизводительной передачи данных между вычислительными узлами, используемая в работе как одна из реализаций интерфейса передачи тензоров}
    \term[Pipeline parallelism (PP)]{метод распределенного выполнения модели, при котором последовательные части модели размещаются на разных устройствах и выполняются как стадии конвейера}
    \term[RDMA (Remote Direct Memory Access)]{механизм прямого доступа к памяти удаленного узла, позволяющий передавать данные с минимальным участием центрального процессора}
    \term[ReduceScatter]{коллективная операция, при которой данные от всех участников сначала объединяются, а затем результат распределяется между участниками по фрагментам}
    \term[Tensor parallelism (TP)]{метод распределенного выполнения модели, при котором отдельные тензорные операции, например матричные умножения, разбиваются между несколькими устройствами}
    \term[Top-k маршрутизация]{способ маршрутизации в MoE слое, при котором для каждого входа выбираются k наиболее подходящих экспертов по оценкам роутера}
    \term[Адресная передача]{передача данных между конкретными участниками распределенной системы без обязательного участия всей коммуникационной группы}
    \term[Асинхронная диспетчеризация]{способ обработки запросов, при котором отправка задач экспертам, ожидание результатов и продолжение работы системы не блокируются одной общей точкой выполнения}
    \term[Авторегрессионный инференс]{режим вывода, при котором модель последовательно предсказывает следующий токен на основе уже имеющегося контекста и ранее полученных токенов}
    \term[Балансировка нагрузки]{распределение запросов или вычислительных задач между доступными экспертными блоками для более равномерной загрузки ресурсов}
    \term[Большая языковая модель]{нейронная модель, обученная на больших объемах текстовых данных и предназначенная для решения задач обработки естественного языка}
    \term[Вычислительный узел]{сервер или отдельный процесс-исполнитель, участвующий в распределенном выполнении модели и располагающий вычислительными ресурсами, например GPU}
    \term[Grouped GeMM (General Matrix Multiplication)]{выполнение нескольких операций матричного умножения как одной сгруппированной операции для уменьшения накладных расходов запуска и более эффективного использования GPU}
    \term[KV кэш]{набор сохраненных ключей и значений механизма внимания, переиспользуемый на последующих шагах авторегрессионного инференса}
    \term[Коллективная операция]{коммуникационная операция, в которой согласованно участвует группа процессов или устройств}
    \term[Коммуникационная группа]{набор участников распределенного вычисления, между которыми выполняются коллективные операции или другие согласованные коммуникации}
    \term[Маршрутизация токенов]{процесс выбора экспертов, которым должны быть отправлены токены или соответствующие им активации в MoE слое}
    \term[Отказоустойчивость]{способность системы сохранять работоспособность или деградировать контролируемым образом при отказе части вычислительных узлов}
    \term[Пропускная способность]{количество запросов, токенов или байтов данных, которое система способна обработать или передать за единицу времени}
    \term[Ранг]{отдельный участник распределенного вычисления, обычно соответствующий процессу, GPU или их связке в рамках коммуникационной группы}
    \term[Роутер]{компонент MoE слоя, определяющий, какие эксперты должны обработать конкретные входные активации}
    \term[Трансформер]{архитектура нейронной сети, основанная на механизме внимания и широко используемая в современных языковых моделях}
    \term[Эксперт]{отдельный вычислительный блок внутри MoE слоя, обычно реализованный как полносвязная нейронная сеть и активируемый только для части входных данных}
    \term[Экспертное вычисление]{этап MoE слоя, на котором выбранные эксперты обрабатывают переданные им активации}
    \term[Экспертный блок]{логически выделенный компонент распределенной MoE системы, который содержит один или несколько экспертов, принимает активации от основной модели, выполняет экспертные вычисления и возвращает результаты; в данной работе экспертный блок также является единицей диспетчеризации, мониторинга и балансировки нагрузки}
    \term[Экспертный слой]{слой MoE модели, содержащий роутер, набор экспертов и операции отправки активаций к экспертам и сборки результатов}
\end{terminologyList}
